% Options for packages loaded elsewhere
\PassOptionsToPackage{unicode}{hyperref}
\PassOptionsToPackage{hyphens}{url}
\PassOptionsToPackage{dvipsnames,svgnames,x11names}{xcolor}
%
\documentclass[
  10pt,
  a4paper,
]{article}
\usepackage{amsmath,amssymb}
\usepackage{setspace}
\usepackage{iftex}
\ifPDFTeX
  \usepackage[T1]{fontenc}
  \usepackage[utf8]{inputenc}
  \usepackage{textcomp} % provide euro and other symbols
\else % if luatex or xetex
  \usepackage{unicode-math} % this also loads fontspec
  \defaultfontfeatures{Scale=MatchLowercase}
  \defaultfontfeatures[\rmfamily]{Ligatures=TeX,Scale=1}
\fi
\usepackage{lmodern}
\ifPDFTeX\else
  % xetex/luatex font selection
  \setmainfont[]{Liberation Serif}
  \setsansfont[]{Liberation Sans}
  \setmonofont[]{DejaVu Sans Mono}
\fi
% Use upquote if available, for straight quotes in verbatim environments
\IfFileExists{upquote.sty}{\usepackage{upquote}}{}
\IfFileExists{microtype.sty}{% use microtype if available
  \usepackage[]{microtype}
  \UseMicrotypeSet[protrusion]{basicmath} % disable protrusion for tt fonts
}{}
\makeatletter
\@ifundefined{KOMAClassName}{% if non-KOMA class
  \IfFileExists{parskip.sty}{%
    \usepackage{parskip}
  }{% else
    \setlength{\parindent}{0pt}
    \setlength{\parskip}{6pt plus 2pt minus 1pt}}
}{% if KOMA class
  \KOMAoptions{parskip=half}}
\makeatother
\usepackage{xcolor}
\usepackage[margin=22mm]{geometry}
\usepackage{longtable,booktabs,array}
\usepackage{calc} % for calculating minipage widths
% Correct order of tables after \paragraph or \subparagraph
\usepackage{etoolbox}
\makeatletter
\patchcmd\longtable{\par}{\if@noskipsec\mbox{}\fi\par}{}{}
\makeatother
% Allow footnotes in longtable head/foot
\IfFileExists{footnotehyper.sty}{\usepackage{footnotehyper}}{\usepackage{footnote}}
\makesavenoteenv{longtable}
\setlength{\emergencystretch}{3em} % prevent overfull lines
\providecommand{\tightlist}{%
  \setlength{\itemsep}{0pt}\setlength{\parskip}{0pt}}
\setcounter{secnumdepth}{5}
\ifLuaTeX
\usepackage[bidi=basic]{babel}
\else
\usepackage[bidi=default]{babel}
\fi
\babelprovide[main,import]{korean}
\ifPDFTeX
\else
\babelfont{rm}[]{Liberation Serif}
\fi
% get rid of language-specific shorthands (see #6817):
\let\LanguageShortHands\languageshorthands
\def\languageshorthands#1{}
\usepackage{xeCJK}
\setCJKmainfont{Noto Serif CJK KR}
\setCJKsansfont{Noto Sans CJK KR}
\setCJKmonofont{Noto Sans Mono CJK KR}
\xeCJKsetup{CJKspace=true}
\usepackage{xurl}
\usepackage{fvextra}
\usepackage{fancyhdr}
\AtBeginDocument{\DefineVerbatimEnvironment{verbatim}{Verbatim}{breaklines=true,breakanywhere=true,fontsize=\small}}
\setlength{\emergencystretch}{4em}
\setlength{\headheight}{15pt}
\pagestyle{fancy}\fancyhf{}
\fancyhead[L]{\small 중성자별 구조와 비방사 진동}
\fancyfoot[C]{\thepage}
\renewcommand{\headrulewidth}{0.3pt}
\setcounter{tocdepth}{2}
\ifLuaTeX
  \usepackage{selnolig}  % disable illegal ligatures
\fi
\usepackage{bookmark}
\IfFileExists{xurl.sty}{\usepackage{xurl}}{} % add URL line breaks if available
\urlstyle{same}
\hypersetup{
  pdftitle={독립 계산과 해석적 검산},
  pdflang={ko-KR},
  colorlinks=true,
  linkcolor={Maroon},
  filecolor={Maroon},
  citecolor={Blue},
  urlcolor={Blue},
  pdfcreator={LaTeX via pandoc}}

\title{독립 계산과 해석적 검산}
\author{}
\date{}

\begin{document}
\maketitle

{
\hypersetup{linkcolor=}
\setcounter{tocdepth}{3}
\tableofcontents
}
\setstretch{1.12}
\section{독립 계산과 해석적
검산}\label{uxb3c5uxb9bd-uxacc4uxc0b0uxacfc-uxd574uxc11duxc801-uxac80uxc0b0}

우리는 같은 물리 문제를 두 실행 경로로 계산했다. 주 계산은 SI 단위의
Fortran TOV·Cowling 식, native EOS 구간 적분과 양방향 DOPRI/RK4
방법이다. 대조한 연구 구현은 기하학적 cm 단위에서 DOP853로 배경을 얻고,
C로 실행하는 GL4 Magnus 전달을 이용해 두 방향의 진동해를 잇는다. 두
경로는 배경 파일이나 적분 결과를 공유하지 않는다.

같은 자료를 서로 다른 보간식에 넣으면 코드 비교에 물리 함수의 차이도
섞인다. 따라서 여기서는 두 경로가 동일한 native cubic \(P(\mu)\)를
평가하도록 독립적인 EOS adapter만 제공했다. 원래 연구 구현의
TOV·전달·matching 코드는 바꾸지 않았다. EOS는 같은 함수이며 적분 방법은
다르다. 이 구분이 비교의 범위다.

\subsection{같은 입력과
단위}\label{uxac19uxc740-uxc785uxb825uxacfc-uxb2e8uxc704}

기본 자료는 84행의 ID69이며 \(P,n,\mu\)의 원래 knot를 유지한다. 각
구간에서 \(n=P'\), \(\epsilon=\mu n-P\), \(\chi=n/(\mu P'')\)를
계산한다. 저밀도 변형은 제10장의 \(\alpha=0.1\)이며 split 행은 69다. 두
별의 목표 질량은 \(2.644606257493485\times10^{30}\) kg, 표면압력은
\(10^{20}\) Pa, 시작 반경은 1 m다. 두 경로 모두 같은 barotropic 응답과
유체 표면조건을 사용한다.

SI 압력과 에너지밀도를 기하학적 cm 단위로 옮기는 계수는
\(G/(10^4c^4)\)다. 반경은 m에서 cm로 100을 곱하고, 연구 코드의 명목
태양질량은 \(1.3271244\times10^{20}/G\) kg으로 환산한다. \(W,V\)의 길이
차원과 공통 진폭 정규화도 맞춘다. 시간 기준은
\(\nu(R)=\tfrac12\ln[1-2GM/(Rc^2)]\)다. 유한 압력 표면이라는 한계는
양쪽에서 동일하다.

연구 코드의 고유값 bracket은 \(f\)에 2000--3000 Hz, \(p_1\)에 6500--8000
Hz를 썼다. 좁은 검색 창에서 얻은 후보에 이름을 자동 부여하지 않고, 두
코드의 고유함수 node와 overlap을 별도로 비교했다. 검색 범위 밖의
스펙트럼까지 완전하다고 주장하지 않는다.

\subsection{배경과 진동 전달을 따로
세분한다}\label{uxbc30uxacbduxacfc-uxc9c4uxb3d9-uxc804uxb2ecuxc744-uxb530uxb85c-uxc138uxbd84uxd55cuxb2e4}

GL4 mesh의 기본 제한은 반경 1000 cm, 로그압력 0.04, 로그반경 0.04다.
아래의 배율은 세 제한을 함께 나누는 수다. 처음 세 수준에서는 배경 DOP853
허용오차를 고정하고 전달 격자만 줄였다. 그 다음 배경 허용오차를 줄여
남은 바닥을 구분했다. 표는 기준 별에서 연구 코드 주파수에서 주 계산
주파수를 뺀 값이다.

\begin{longtable}[]{@{}
  >{\raggedleft\arraybackslash}p{(\columnwidth - 6\tabcolsep) * \real{0.2500}}
  >{\raggedleft\arraybackslash}p{(\columnwidth - 6\tabcolsep) * \real{0.2500}}
  >{\raggedleft\arraybackslash}p{(\columnwidth - 6\tabcolsep) * \real{0.2500}}
  >{\raggedleft\arraybackslash}p{(\columnwidth - 6\tabcolsep) * \real{0.2500}}@{}}
\toprule\noalign{}
\begin{minipage}[b]{\linewidth}\raggedleft
격자 배율
\end{minipage} & \begin{minipage}[b]{\linewidth}\raggedleft
배경 상대 허용오차
\end{minipage} & \begin{minipage}[b]{\linewidth}\raggedleft
\(f\) 차이 {[}Hz{]}
\end{minipage} & \begin{minipage}[b]{\linewidth}\raggedleft
\(p_1\) 차이 {[}Hz{]}
\end{minipage} \\
\midrule\noalign{}
\endhead
\bottomrule\noalign{}
\endlastfoot
1 & \(2\times10^{-11}\) & \(-1.098\times10^{-8}\) &
\(-6.282\times10^{-7}\) \\
2 & \(2\times10^{-11}\) & \(-4.554\times10^{-9}\) &
\(-4.384\times10^{-8}\) \\
4 & \(2\times10^{-11}\) & \(-4.159\times10^{-9}\) &
\(-7.537\times10^{-9}\) \\
4 & \(2\times10^{-13}\) & \(-2.055\times10^{-10}\) &
\(-2.766\times10^{-9}\) \\
8 & \(2\times10^{-13}\) & \(-1.810\times10^{-10}\) &
\(-5.030\times10^{-10}\) \\
\end{longtable}

차이가 작아졌다는 사실과 개별 코드의 오차가 그 차이보다 작다는 주장은
같지 않다.{[}A4{]} 같은 EOS와 가정을 공유하는 오류가 남을 수 있으므로,
별도의 식·경계·구적 검사를 함께 사용한다.

\subsection{기준 별과 변형
별}\label{uxae30uxc900-uxbcc4uxacfc-uxbcc0uxd615-uxbcc4}

\begin{longtable}[]{@{}lrr@{}}
\toprule\noalign{}
모드 & 주 계산 기준 별 {[}Hz{]} & 연구 계산 기준 별 {[}Hz{]} \\
\midrule\noalign{}
\endhead
\bottomrule\noalign{}
\endlastfoot
\(f\) & 2521.359983099551 & 2521.359983099370 \\
\(p_1\) & 7253.656215854635 & 7253.656215854132 \\
\end{longtable}

\begin{longtable}[]{@{}lrr@{}}
\toprule\noalign{}
모드 & 주 계산 변형 별 {[}Hz{]} & 연구 계산 변형 별 {[}Hz{]} \\
\midrule\noalign{}
\endhead
\bottomrule\noalign{}
\endlastfoot
\(f\) & 2521.349534625408 & 2521.349534625234 \\
\(p_1\) & 7253.186474083501 & 7253.186474083005 \\
\end{longtable}

두 코드의 차이는 기준·변형 모두 \(f\)에서 약 \(2\times10^{-10}\) Hz,
\(p_1\)에서 약 \(5\times10^{-10}\) Hz다. 변형 별의 연결점을 0.5에서
0.65의 위치 분율로 바꾼 차이도 각각 \(2\times10^{-12}\) Hz 미만이었다.
두 모드의 node는 각각 \((0,0),(1,1)\)이며 아홉 끝점·영점 기준에서
유지되었다. 공통 좌표·정규화의 squared overlap은 모든 대조에서
\(0.9999999998\)보다 컸다. 수치 비교의 자릿수는 입력 자료의 물리적
정확도와 다르다.

연구 구현에 저장된 별도의 분석적 폴리트로프·물질적 계면·고정질량 사례도
소스에서 다시 계산했다. 그 사례의 \(f=2089.581136931826\) Hz,
\(p_1=4597.638996128369\) Hz와 다섯 후보 주파수가 저장 정밀도에서
재현되었다. 이는 저장된 실행의 재현성 검사이며 위 연속 ID69 문제와 같은
별이라는 뜻은 아니다. 물질적 계면 스펙트럼 전체를 주 계산과 교차검증한
결과로 확대하지 않는다.

\subsection{고유함수로부터 독립적으로 주파수를
검사한다}\label{uxace0uxc720uxd568uxc218uxb85cuxbd80uxd130-uxb3c5uxb9bduxc801uxc73cuxb85c-uxc8fcuxd30cuxc218uxb97c-uxac80uxc0acuxd55cuxb2e4}

제5장에서 유도한 scalar Rayleigh 식에는 유한 표면의 양의 분모 항과 유한
시작점의 flux가 있다. 주 계산이 반환한 \(W,V\)에서 주파수를 되넣지 않고
국소 시험함수를 구성하여 적분했다. 가장 정밀한 기준 별에서 \(f\)의
적분값은 2521.359983089991 Hz, \(p_1\)은 7253.656215853673 Hz였다.
Shooting 결과와의 차이는 각각 \(-9.56\times10^{-9}\) Hz,
\(-9.62\times10^{-10}\) Hz다. 변형 별에서도 같은 규모의 일치를 얻었다.

분모에서 표면항을 잘못 생략하면 같은 시험함수의 재구성 주파수는
\(f\)에서 약 \(1.07\times10^{-6}\) Hz, \(p_1\)에서 약
\(5.60\times10^{-5}\) Hz만큼 바뀐다. 실제 solver가 이 항을 빠뜨렸다는
뜻은 아니다. Solver는 이미 올바른 표면조건을 사용한다. 잘못된 검산식을
사용하면 정상적인 해를 부정확하다고 오판할 수 있다는 반례다.

4점·8점 구적의 차이와 profile refinement는 구분했다. 구적 차이가 0으로
반올림되더라도 고유함수·배경 재구성 오차의 상한은 얻어지지 않는다. 최종
해석에는 독립 코드 비교, 절대 주파수와 paired 차이의 수렴, 모드 대응,
중심·표면·질량·연결점 민감도를 함께 사용한다.

\subsection{중심에서의 분석적
개선}\label{uxc911uxc2ecuxc5d0uxc11cuxc758-uxbd84uxc11duxc801-uxac1cuxc120}

같은 \(P_c=1.5\times10^{34}\) Pa에서 시작 반경을 1 m에서 100 m로 바꾸어
보았다. 이 비교에서는 질량을 다시 맞추지 않고 중심 초기값의 영향을
측정했다.

\begin{longtable}[]{@{}
  >{\raggedright\arraybackslash}p{(\columnwidth - 4\tabcolsep) * \real{0.3333}}
  >{\raggedleft\arraybackslash}p{(\columnwidth - 4\tabcolsep) * \real{0.3333}}
  >{\raggedleft\arraybackslash}p{(\columnwidth - 4\tabcolsep) * \real{0.3333}}@{}}
\toprule\noalign{}
\begin{minipage}[b]{\linewidth}\raggedright
시작법
\end{minipage} & \begin{minipage}[b]{\linewidth}\raggedleft
\(f\)의 1 m 기준 차이 {[}Hz{]}
\end{minipage} & \begin{minipage}[b]{\linewidth}\raggedleft
\(p_1\)의 1 m 기준 차이 {[}Hz{]}
\end{minipage} \\
\midrule\noalign{}
\endhead
\bottomrule\noalign{}
\endlastfoot
100 m, leading 시작 & \(-5.51\times10^{-6}\) & \(-1.05\times10^{-5}\) \\
100 m, 다음 차수 포함 & \(+7.09\times10^{-11}\) &
\(+1.60\times10^{-10}\) \\
\end{longtable}

허용오차는 모두 \(10^{-12}\)로 같았다. 배경의 accepted step은 3037에서
2307로, \(f\) 적분은 3063에서 2333으로 줄었다. 이것은 측정한 한 모형의
결과다. 모든 EOS나 고차 모드에서 같은 개선 배율이나 실행시간 단축을
보장하지 않는다. 중심 급수는 표준 물리식을 사용할 때만 선택하며, 기본
비교값은 계속 1 m의 leading 시작을 사용한다.

\subsection{결론의 범위}\label{uxacb0uxb860uxc758-uxbc94uxc704}

현재의 연속 ID69 모형과 유한 표면의 두 모드는 서로 다른 실행 경로와 독립
적분 검사를 통과했다. 같은 질량의 변형 차이도 검토한 수치 여유보다 크다.
이것은 주어진 Cowling·barotropic 모형 안의 수치적 신뢰를 지지한다. 실제
조성 반응, 탄성, 회전, 자기장, 초유체나 metric perturbation의 불확실성을
줄인 결과는 아니다.{[}C3{]}

선택된 연구 코드 검사는 18개의 Python 시험과 네 Fortran 프로그램의 82개
assertion을 포함했다. 전체 저장 사례와 모든 parameter를 새로 실행한 것은
아니다. 학생용 기본 계산에는 외부 연구 저장소나 SciPy가 필요하지 않고,
별도 코드 대조를 직접 재현할 때만 지정한 연구 저장소와 SciPy를 사용한다.

\subsection{읽기}\label{uxc77duxae30}

{[}A4{]} R. H. D. Townsend, R. V. Kuenzi and J. Christensen-Dalsgaard,
``Tools for Characterizing the Numerical Error of Stellar Oscillation
Codes,'' arXiv:2505.20212 (2025). 코드 비교와 개별 해의 오차 특성화를
구분한다.

{[}C3{]} G. Montefusco, M. Antonelli and F. Gulminelli, ``Frozen and
\(\beta\)-equilibrated \(f\) and \(p\) modes of cold neutron stars:
nuclear metamodel predictions,'' \emph{A\&A} \textbf{694}, A150 (2025),
arXiv:2410.08008. 평형과 고정 조성 압축응답의 물리적 차이를 다룬다.

\end{document}
